\clearpage
\section*{Problem 9}
\addcontentsline{toc}{section}{Problem 9}
\markboth{Problem 9}{Problem 9}

\textbf{Problem Statement (Textbook 6.1.8):} 

Discuss the problem of determining a polynomial of degree at most 2 for which $p(0)$, $p(1)$, and $p(\xi)$ are prescribed, $\xi$ being any preassigned point.

\noindent\rule{\textwidth}{0.4pt}

\vspace{1em}

\textbf{Solution:}

We seek a polynomial $p(x) = a_0 + a_1 x + a_2 x^2$ (degree at most 2) such that:
$$p(0) = y_0, \quad p(1) = y_1, \quad p(\xi) = y_\xi$$
where $y_0, y_1, y_\xi$ are prescribed values and $\xi$ is a given point.

\subsection*{Main Case: Three Distinct Points Required}

For a polynomial of degree at most 2, we need three distinct interpolation points. Thus, we require:
$$\boxed{\xi \neq 0 \text{ and } \xi \neq 1}$$

The Vandermonde system is:
$$\begin{bmatrix} 
1 & 0 & 0 \\
1 & 1 & 1 \\
1 & \xi & \xi^2
\end{bmatrix}
\begin{bmatrix} a_0 \\ a_1 \\ a_2 \end{bmatrix}
=
\begin{bmatrix} y_0 \\ y_1 \\ y_\xi \end{bmatrix}$$

The determinant is:
$$\det(V) = \xi(\xi - 1)$$

Since $\xi \neq 0$ and $\xi \neq 1$, we have $\det(V) \neq 0$, so the system has a \textbf{unique solution}.

\textbf{Conclusion:} When $\xi \notin \{0, 1\}$, there exists a \textbf{unique polynomial of degree at most 2} satisfying all three conditions.

\subsection*{Degenerate Cases}

\textbf{Note:} If $\xi = 0$ or $\xi = 1$, we only have two distinct interpolation points, which contradicts the requirement for a degree-2 polynomial. In these cases:
\begin{itemize}
\item The problem reduces to linear interpolation (degree at most 1)
\item The prescribed values must be consistent (e.g., if $\xi = 0$, then $y_0 = y_\xi$)
\item The resulting polynomial is the line through the two distinct points
\end{itemize}

These cases defeat the original question's intent of determining a degree-2 polynomial.

\textbf{Key Insight:} The problem is well-posed if and only if we have three distinct interpolation points, i.e., $\xi \notin \{0, 1\}$.

\vspace{2em}
\noindent\rule{\textwidth}{0.4pt}
